\documentclass{article}

% For camera-ready workshop papers, use:
\usepackage[dblblindworkshop, final]{neurips_infpriv_2026}

% To compile a preprint version, e.g., for submission to arXiv:
% \usepackage[preprint]{neurips_infpriv_2026}

% To compile a double-blind workshop submission:
% \usepackage[dblblindworkshop]{neurips_infpriv_2026}

\usepackage[utf8]{inputenc} % allow utf-8 input
\usepackage[T1]{fontenc}    % use 8-bit T1 fonts
\usepackage{hyperref}       % hyperlinks
\usepackage{url}            % simple URL typesetting
\usepackage{booktabs}       % professional-quality tables
\usepackage{amsfonts}       % blackboard math symbols
\usepackage{nicefrac}       % compact symbols for 1/2, etc.
\usepackage{microtype}      % microtypography
\usepackage{xcolor}         % colors

\title{Your Paper Title Here}

\author{%
  Author One \\
  Department of Computer Science \\
  Institution \\
  \texttt{author1@institution.edu} \\
  \And
  Author Two \\
  Department of Computer Science \\
  Institution \\
  \texttt{author2@institution.edu} \\
}

\begin{document}

\maketitle

\begin{abstract}
The camera-ready version for InfPriv main track is limited to 5 pages, excluding references and appendices. The camera-ready version for InfPriv fast track is limited to 10 pages, excluding references and appendices. The camera-ready version of the fast-track paper can remove its appended review. NeurIPS checklist is optional. All camera-ready version should use the ``neurips\_infpriv\_2026'' style file with ``dblblindworkshop, final'' options. All other formatting requirements follow the NeurIPS 2026 guidelines.
\end{abstract}

\section{Introduction}
Provide the introduction to your research here.

\section{Methodology}
Detail your method and theoretical or experimental formulation.

\section{Experiments and Results}
Present your empirical evaluation and findings.

\section{Conclusion}
Summarize contributions and future directions.

\bibliographystyle{plain}
% \bibliography{references}

\end{document}
